\scanpage{192}
4) Soit $X$ de type fini sur un corps, $X$ \deleted{\unsure{géom.}} \uuline{géom.} \uuline{unibranche}, on veut étudier les modules inversibles sur $X$ à stratification intégrable, \deleted{\unsure{indécomposable|de descente|nécessaire}}. \margin{loc. de t.f. suffit ?}

Soit $\underline L$ un module \inserted{inversible} stratifié sur $X$. Je dis que la stratification est intégrable ssi $\exists$ entier $n>0$, \inserted{premier à la car.,} tel que $\underline L^{\otimes n}$ soit trivial comme module stratifié [Le ``si'' est valable sans hyp. géom. unibranche], i.e ssi \deleted{on trouve} \inserted{$\underline L$ donne} un élément d'ordre fini \inserted{premier à $p$} dans le groupe des modules inversibles stratifiés. A fortiori, une \inserted{première} condition nécessaire est qu'il existe un entier $n>0$ \inserted{premier à $p$} tel que $\underline L^{\otimes n}$ admette une section partout $\ne 0$, i.e. que $\underline L$ donne élément d'ordre fini \inserted{premier à $p$} dans le $\operatorname{Pic}(X)$.

Cette condition étant remplie, il reste à exprimer que $\underline L^{\otimes n}$ est à \deleted{connexion} \inserted{stratification} \uline{intégrable}. On est donc ramené à caractériser \deleted{\unsure{les}} parmi les \deleted{connexions} \inserted{stratifications} sur $\underline O_X$, \deleted{les connexions} \inserted{celles qui sont} intégrables. Or, si $X/k$ est lisse, \deleted{\unsure{et $k$ de car. 0|et $k$ de car. nulle}} \deleted{\unsure{sur $k$ de car. nulle}} les stratifications sur $\underline O_X$ correspondent aux \inserted{connexions à courbure nulle, i.e. aux} 1-formes \uline{fermées} $\omega$ sur $X$. \margin{\uline{ici $k$ de car. 0} !} Dire que la stratification \inserted{(ou la connexion)} est triviale \deleted{\unsure{(ou encore, la conn.|(ou connex., la conn.}} revient à dire que \uline{$\omega$ est une dérivée logarithmique}. Dire qu'elle est intégrable revient donc à dire qu'il y a un entier $n>0$ tel que $n\omega$ soit une dérivée logarithmique. Comme le groupe des dérivées logarithmiques est un groupe isomorphe à un $L=\mathbb Z^{\unsure{r|n}}$ [\unsure{Rappeler|Rappel de|d’après} Rosenlicht], il est clair
% archive sheet number: 191
